


\section{Causal generators and feasible policy transfer}
Acciaio, Eckstein and Hou construct a time-causal variational autoencoder whose
encoder and decoder at date $t$ use only input coordinates available by $t$.
Their retained preprint bounds a \emph{directional} causal transport discrepancy
by reconstruction error and a term involving latent relative entropy
\cite[Lemma~2.4, Theorem~2.5 and Corollary~2.6]{litclosetcvae}. The relative-entropy
step requires the reconstructed and generated distributions to have the stated
bounded support. Normalizing the observed sample does not establish that a
learned decoder satisfies that support restriction for every generated input.
A finite training estimate is also different from a population bound.

There is a further issue in transferring the control result. Theorem~3.1 of
the retained version states its comparison for a closed convex subset
$\mathcal H$ of adapted strategies. Its proof takes a policy $G$ on the second
process and constructs $\widetilde G(x)=\E[G(Y)\mid X=x]$. Causality establishes
adaptedness, but adaptedness alone does not establish
$\widetilde G\in\mathcal H$. Here is a direct counterexample to that stated
generality. There is one date, $\mu=\delta_1$, $\nu=\delta_0$, and
\begin{equation}
\mathcal H=\{h\},\qquad h(x)=\min(1,\max(0,x)),\qquad Q(x,a)=a.
\end{equation}
The singleton class is closed and convex, and $Q$ is convex in $a$ and
uniformly zero-Lipschitz in $x$. Nevertheless,
$V(\mu)-V(\nu)=1$, whereas the proposed right-hand side with $L=0$ is zero.
The transferred action at $x=1$ is zero, outside the singleton policy class.
This disproves the unrestricted policy-class statement in that retained
version. It does not refute the argument when the feasible class is explicitly
closed under the conditional transfer. The published edition has not been
matched equation by equation, so this is not asserted to be a journal erratum.

For AccDSL, state-dependent borrowing limits, integer issuance sizes and
contractual exercise restrictions make closure substantive. A repaired transfer
argument must verify feasibility of the mapped policy, or use a method that
retains the original feasible actions. The one-sided value comparison is not
itself a bound on the performance of the same learned policy after deployment.
The inspected paper-linked code has forward LSTM encoder/decoder paths, but
also external conditioning variables; those variables require their own
availability dates. Its selected conditional decoder uses ReLU and a final
linear layer, so the architecture does not enforce the theorem's unit-ball
support. These observations identify additional obligations before adoption;
no training result is inferred from source-code inspection.

\section{Controlled transition laws and off-policy support}
Sald\i{} and Y\"uksel distinguish weak and strong kernel-embedding topologies
and give a square-wave counterexample to their equivalence
\cite[Theorem~4.1 and proof]{litclosekerneltop}. Let $Y$ be uniform on $[0,1]$.
On alternating half-cells of a fine partition, make a binary successor $U$
deterministically 1 or 0. As the partition becomes finer, smooth integrated
tests can approach those of a fair binary kernel. At every particular input,
however, the conditional law remains a point mass. For the Kronecker kernel
on $\{0,1\}$, its squared embedding discrepancy from the fair law is $1/2$.
The local exact calculation checks this discrepancy and the vanishing error
for a linear test, rather than claiming that a few tests prove weak convergence.

Their discounted-control application uses bounded continuous costs, finite
actions, continuous transition kernels and a common dominating measure, including
for the initial distribution. One route requires a strong essential-supremum
condition; the weaker route adds bounded equicontinuous densities
\cite[Assumptions~2--4, equation~(11), Theorems~5.1--5.2]{litclosekerneltop}.
A default discontinuity, hidden creditor information or an action-dependent
financing constraint therefore needs a separate argument. An equally important
problem is support: models that agree under the historically chosen action may
disagree under a new action. With action losses $(1,100)$ in the real model and
$(1,0)$ in an approximate model, historical use of the first action gives
identical observations. Optimizing the approximation chooses the second action
and incurs regret 99. Agreement under the historical policy alone cannot
establish a policy-uniform approximation claim.

The original robustness result used in that application makes another condition
explicit. Kara and Y\"uksel distinguish continuity of optimal values from
robustness of the policies attaining them. Their discounted-control argument
uses continuous kernel convergence, bounded continuous costs and compact
actions. Policy transfer additionally requires each approximate policy to be
optimal from every initial state, as when it satisfies the discounted Bellman
equation throughout the state space
\cite[Assumption~1.1, Theorems~4.2--4.4 and Appendix~A.5]{litclosekara}.
An optimizer trained only from a historical starting state can leave its actions
elsewhere arbitrary; the true transition can reach those neglected states.
The paper's counterexample makes precisely this distinction. Its finite-state
learning result also assumes continuing coverage of every state--action pair
\cite[Corollary~5.1]{litclosekara}. Financial statements generated by one management
policy do not by themselves provide that coverage.
